\documentclass{article}
\usepackage[UTF8]{ctex}
\usepackage{amsmath,amssymb,amsfonts}
\usepackage{enumitem}
\usepackage{geometry}
\geometry{a4paper, margin=1in}
% === 错题本排版优化（不参与题目标识） ===
\setlength{\parindent}{0pt}
\setlength{\parskip}{0.35em plus 0.10em minus 0.08em}
\setlength{\emergencystretch}{3em}
\setlist{itemsep=0.25em, topsep=0.40em, parsep=0pt}
\allowdisplaybreaks[2]
% === 错题本排版优化结束 ===
\begin{document}
\section*{西南大学2025年高等代数真题}
\section*{57．西南大学2025年高等代数试卷}

\subsection{题目1}
1．（15 分）令 \(m, n\) 为正整数，且 \(m>n, f(x)\) 是有理数域上的 \(n\) 次多项式，\(p\) 为素数，证明：\(f(\sqrt[m]{p}) \neq 0\)

\subsection{题目2}
2．（15 分）计算 \(n(n \geq 2)\) 阶行列式 \(\left|s_{i j}\right|\) ，其中 \(s_{i j}=a^{i+j-2}+b^{i+j-2}\) 。

\subsection{题目3}
3．（15 分）令 \(F\) 为数域，\(A \in F^{n \times m}, B \in F^{n \times n}\) ，证明：矩阵方程组 \(A X=B\) 有解当且仅当 \(A\) 的秩等于分块矩阵 \((A, B)\) 的秩．

\subsection{题目4}
4．（15 分）令 \(F\) 为数域，问矩阵空间 \(F^{n \times n}\) 是否存在由可逆矩阵形成的基底？证明你的结论。

\subsection{题目5}
5．（20 分）已知

\[
\left\{\begin{array}{l}
x_{1}+x_{2}+x_{3}=1 \\
x_{1}+2 x_{2}+3 x_{3}=2 \\
a x_{1}+3 x_{2}+b x_{3}=3
\end{array}\right.
\]

有两个线性无关的解，求 \(a, b\) 的值，并求该方程组的一般解．

\subsection{题目6}
6．（20 分）已知 3 阶矩阵 \(A\) 满足 \(|A-E|=|A-2 E|=|A+E|=a\) ．

\begin{itemize}
  \item[（1）] 当 \(a=0\) 时，求行列式 \(|A+3 E|\) 。
  \item[（2）] 当 \(a=2\) 时，求行列式 \(|A+3 E|\) 。
\end{itemize}

\subsection{题目7}
7．（20 分）令 \(F\) 为数域，\(V=F^{n \times n}, A \in V\) 。定义 \(V\) 上的线性变换 \(\sigma\) 如下：

\[
\sigma(X)=A X-X A .
\]

（1）证明：\(\sigma\) 不可逆．\\
（2）设 \(n=2, A=\left(\begin{array}{ll}3 & -4 \\ 2 & -3\end{array}\right)\) ，判断 \(\sigma\) 是否可对角化？并给出理由．

\subsection{题目8}
8．（15 分）令 \(\sigma, \tau\) 是 \(n\) 维欧氏空间 \(V\) 上的线性变换，且对任意的 \(\alpha\) ，都有 \((\sigma(\alpha), \sigma(\alpha))=(\tau(\alpha), \tau(\alpha))\)－证明：存在正交变换 \(\xi\) ，使得 \(\tau=\xi \sigma\) ．

\subsection{题目9}
9．（15 分）令 \(A, B\) 均为 \(n\) 阶正定矩阵，证明：\(A B\) 正定当且仅当 \(A B=B A\) 。
\end{document}
